\begin{tikzpicture}[
  >={Latex[length=3mm,width=2.2mm]},
  every node/.style={font=\small},
  bvec/.style={->,line width=1.3pt,bblue},
  dashline/.style={dashed,thin,gray!60!black}
]

% Asse orizzontale
\draw[->,thick] (-5,0) -- (5,0);

% Filo infinito (cerchietti viola con punto = corrente uscente)
\foreach \x in {-4.5,-4.0,...,-0.5}{
  \draw[vpurple!40,fill=vpurple!15] (\x,1.2) circle (0.2);
  \fill[vpurple] (\x,1.2) circle (1pt);
}

% dx' indicazione
\draw[<->,thin] (-2.5,1.6) -- (-2.0,1.6);
\node[above] at (-2.25,1.6) {$dx'$};

% R verticale tratteggiata
\draw[dashline] (-2.25,0) -- (-2.25,1.0);
\node[left] at (-2.25,0.5) {$R$};

% r diagonale
\draw[thin] (-2.25,1.2) -- (2.5,0);
\node[above] at (0.5,0.6) {$r$};

% Angolo phi
\draw[thin] (1.5,0) arc (180:155:1.0);
\node[above left] at (1.5,0.15) {$\phi$};

% Punto P
\coordinate (P) at (2.5,0);
\fill (P) circle (2pt);
\node[above] at (P) {$P$};

% dB
\draw[bvec] (P) -- ++(0.8,0) node[right] {$d\mathbf{B}$};

% Etichette
\node[below] at (-2.25,-0.15) {$x'$};
\node[below] at (2.5,-0.15) {$x$};

% --- In basso: filo con corrente entrante (x = croce) ---
\begin{scope}[yshift=-1.8cm]
  \foreach \x in {-4.5,-4.0,...,4.0}{
    \draw[vpurple!40,fill=vpurple!10] (\x,0) circle (0.2);
    \draw[vpurple,thin] ({\x-0.1},-0.1) -- ({\x+0.1},0.1);
    \draw[vpurple,thin] ({\x-0.1},0.1) -- ({\x+0.1},-0.1);
  }
  \draw[<->,thin] (-2.0,-0.5) -- (2.0,-0.5);
  \node[below] at (0,-0.5) {$d$};
\end{scope}

\end{tikzpicture}
